\par\Needspace{6\baselineskip}\medskip
\hypertarget{native-208-article154}{}\pdfbookmark[2]{154.}{article-154}
\begin{center}154.\end{center}
To leave nothing wanting, we shall construct a \textit{third hypothesis}, again adopting the values of $P',Q'$ obtained in the second hypothesis as the values of $P,Q$. Thus, setting
\[
\begin{aligned}
x&=\log P=0{,}079\,0167\\
y&=\log Q=8{,}547\,6110,
\end{aligned}
\]
we obtain the following principal results of the computation:
\begin{center}\small\setlength{\tabcolsep}{6pt}
\begin{tabular}{@{}r|r@{}}
$\omega=13^{0}\,15'\,38''{,}39$ & $\zeta''=210^{0}\,8'\,25''{,}65$\\[.5em]
$\omega+\sigma=13\quad38\quad51{,}51$ & $\log r=0{,}330\,7640$\\[.5em]
$\log Qc\sin\omega=0{,}598\,9542$ & $\log r''=0{,}322\,2239$\\[.5em]
$z=14^{0}\,33'\,19''{,}50$ & $\tfrac12(u''+u)=205^{0}\,22'\,14''{,}57$\\[.5em]
$\log r'=0{,}325\,9878$ & $\tfrac12(u''-u)=-3\quad14\quad4{,}78$\\[.5em]
$\log\dfrac{n'r'}{n}=0{,}667\,5154$ & $2f'=7\quad34\quad53{,}73$\\[.7em]
$\log\dfrac{n'r'}{n''}=0{,}588\,4987$ & $2f=3\quad29\quad0{,}39$\\[.7em]
$\zeta=203^{0}\,16'\,38''{,}41$ & $2f''=4\quad5\quad53{,}34$
\end{tabular}
\end{center}

All these numbers differ so little from those furnished by the second hypothesis that we may safely conclude that the third hypothesis requires no further correction\footnote{If the computation were carried through to the end as in the preceding hypotheses, it would give $X=0$, $Y=-0{,}000\,0003$. The latter value may be regarded as vanishing and scarcely exceeds the uncertainty that always attaches to the last decimal digit.}. We may therefore proceed to determine the elements themselves from $2f',r,r'',\theta'$. We omit the details here, since this determination has already been given in full as an example in Article~97. It remains only to compute the position of the orbital plane according to Article~149 and transfer the epoch to the beginning of 1805. This computation is to be based on the following numbers:
\SourceLeaf{209}
\[
\begin{aligned}
AD'-\zeta&=9^{0}\,55'\,51''{,}41\\
\tfrac12(\gamma+u)&=202\quad18\quad13{,}855\\
\tfrac12(\gamma-u)&=-6\quad18\quad5{,}495,
\end{aligned}
\]
from which we derive
\[
\begin{aligned}
\tfrac12(g+h)&=196^{0}\,43'\,14''{,}62\\
\tfrac12(g-h)&=-4\quad37\quad24{,}41\\
\tfrac12i&=6\quad33\quad22{,}05.
\end{aligned}
\]
Thus $h=201^{0}\,20'\,39''{,}03$, and hence $\AscendingNode=l-h=171^{0}\,7'\,48''{,}73$. Further, $g=192^{0}\,5'\,50''{,}21$. Since the true anomaly for the first observation was found in Article~97 to be $=310^{0}\,55'\,29''{,}64$, the distance of the perihelion from the ascending node in the orbit is $=241^{0}\,10'\,20''{,}57$, and the longitude of the perihelion $=52^{0}\,18'\,9''{,}30$. Finally, the inclination of the orbit is $=13^{0}\,6'\,44''{,}10$. --- If we prefer to start the same computation from the third position, we have
\[
\begin{aligned}
A''D'-\zeta''&=24^{0}\,18'\,35''{,}25\\
\tfrac12(\gamma''+u'')&=196\quad24\quad54{,}98\\
\tfrac12(\gamma''-u'')&=-5\quad43\quad14{,}81.
\end{aligned}
\]
This gives
\[
\begin{aligned}
\tfrac12(g''+h'')&=211^{0}\,24'\,32''{,}44\\
\tfrac12(g''-h'')&=-11\quad43\quad48{,}49\\
\tfrac12i&=6\quad33\quad22{,}05,
\end{aligned}
\]
and hence the longitude of the ascending node $=l''-h''=171^{0}\,7'\,48''{,}72$, the longitude of the perihelion $=52^{0}\,18'\,9''{,}30$, and the inclination of the orbit $=13^{0}\,6'\,44''{,}10$, all just as before.

The interval from the last observation to the beginning of 1805 is $64{,}614\,102$ days. The corresponding mean heliocentric motion is $53293''{,}65=14^{0}\,48'\,13''{,}65$. Thus the epoch of mean anomaly for the beginning of 1805 on the Paris meridian is $=349^{0}\,34'\,12''{,}38$, and the epoch of mean longitude $=41^{0}\,52'\,21''{,}68$.
